\subsection{One count for the two initial energy classes}

The compact and source functions describe the same return on their
overlap. We identify their positive multiplier and show that their
divided coefficient triples have the same limit there. This permits
us to join the coefficients while preserving the Wronskian signs and
count the two energy classes together.

\begin{proposition}[The initial joint bound]
\label{rw:initial-joint-bound}
Keep the collar and the base threshold of Section~\ref{rw:geometry}.
There are constants $\eta>0$ and $H_*>h(z_L)+1$, followed by a
reflection-invariant original coefficient neighborhood, such that the
two initial classes in \eqref{rw:geometry:energy-classes} together
contain at most two isolated counted cycles, including multiplicity,
for each fixed field. If $\ell=0$, their joint bound is one, and
that cycle is simple. If also $\tau_0=\sigma_0=0$, the corresponding
return is the identity and contributes no isolated cycle.
\end{proposition}

We use the same lower-section energy as in Appendix
\ref{rw:appendix-base-ect}:
\[
 h(s)=s-1-\log s,\qquad
 \Delta(h,\alpha)=h(P(h,\alpha))-h.
\]
Here $P$ is the actual positive central return, analytic on the compact
energy tube used in \eqref{rw:ect-coefficients}. The source construction
uses the fixed-field direction $\ell\ge0$; reflection transfers the
count at periodic points afterwards.

\begin{lemma}[Compact coefficient division]
\label{rw:source:compact-division}
Fix $[h_1,h_2]\Subset(H_{\rm ECT},\infty)$ and a slightly larger
compact interval in that tail. On a sufficiently small coefficient
neighborhood the analytic positive return on this larger interval has
the exact representation
\begin{equation}\label{rw:source:compact-columns}
 \Delta=\tau_0\mathcal A+\sigma_0\mathcal C+\ell\mathcal B.
\end{equation}
The coefficient triple $(\mathcal A,\mathcal C,\mathcal B)$
converges in $C^2([h_1,h_2])$ as the coefficients tend to zero to
$(\mathfrak a,\mathfrak c,\mathfrak b)$ from Appendix
\ref{rw:appendix-base-ect}. Consequently its three ordered Wronskians
have signs $+,+,-$ on this fixed interval.
\end{lemma}

\begin{proof}
Use independent coefficient coordinates
$(\tau_0,\ell,a,t,B)$, where $t=B+m$ and $c=\tau_0-a$.
First divide in the line-breaking coefficient:
\[
 \Delta=\Delta_0+\ell\mathcal B,\qquad
 \Delta_0=\Delta|_{\ell=0},\qquad
 \mathcal B=\int_0^1\partial_\ell\Delta
                 (h,\tau_0,v\ell,a,t,B)\,\dd v.
\]
All coefficient segments stay in the common analytic return tube.
On $\tau_0=a=\ell=0$ the field is reversible; on
$\tau_0=t=\ell=0$ it is on the second center component.
After a coefficient shrink, the real period domains of
Section~\ref{rw:centers} contain the fixed compact family of base
section points on both faces. The positive local return fixes these
periodic inputs, since an increasing interval map has no nontrivial
finite periodic orbit.

Apply the same specified integral rule
\eqref{rw:source:Hadamard-rule}, now to
$\Psi(T,u,t,B)=\Delta_0(h,T,u,t,B)$ with $h$ passive.
It gives $\Delta_0=\tau_0M+atN$. Since
$at=\tau_0t-\sigma_0$, take
$\mathcal A=M+tN$ and $\mathcal C=-N$.
At zero parameter these are respectively
$\partial_{\tau_0}\Delta$, $-\partial_a\partial_t\Delta$
and $\partial_\ell\Delta$, with precisely the normalization
in \eqref{rw:ect-coefficients}. Analyticity on a neighborhood of the
fixed compact product gives their $C^2$ convergence.
Corollary~\ref{rw:base-ect-compact} preserves the strict Wronskian
signs after a further coefficient shrink.
\end{proof}

\begin{lemma}[Identification with the source displacement]
\label{rw:source:common-displacement}
On the positive-core source interval in the direction $\ell\ge0$,
let $\mathcal D$ be the exact source function in
\eqref{rw:source:matched}; thus
\begin{equation}\label{rw:source:D-columns}
 \mathcal D=\tau_0 f+\sigma_0 qf-\Phi,
 \qquad f=h^{3/2}A>0,\quad q=hC/A.
\end{equation}
Then
\begin{equation}\label{rw:source:positive-unit}
 \mathcal D=W(h,\alpha)\Delta(h,\alpha),\qquad W>0.
\end{equation}
The factor is locally analytic on the actual positive-return domain.
On each fixed compact overlap with the analytic return tube, the
source coefficient triple
\[
 (f,qf,-\Phi/\ell)
\]
has a removable third member at $\ell=0$ and converges in $C^2$ to
\begin{equation}\label{rw:source:common-base-triple}
 W(h,0)(\mathfrak a,\mathfrak c,\mathfrak b).
\end{equation}
The comparison is made after dividing that member by $\ell$.
\end{lemma}

\begin{proof}
Fix the source label $h$ and its actual field. Regard the core section
$|x|=\sqrt h$ as fixed while changing a central comparison height.
Let $E_h(u)$ be the incoming action at $y=0$ obtained by following
the incoming comparison path backwards from central height $u$.
The outgoing core, tail and upper arc give an action $E_{\mathrm{out}}$.
For the actual positive return $P$,
\[
 E_{\mathrm{out}}=E_h(P),\qquad E_{\mathrm{in}}=E_h(s).
\]
Scalar order keeps the core solutions, tail inputs and actions between
their positive endpoint paths in the common inverse intervals of
Proposition~\ref{rw:source-maps}. The divided differences below
therefore use actual paths.
Scalar variation gives $z_-'(u)>0$ for the incoming core endpoint.
Its logarithmic entry $-h^{-1}\log z_-(u)$ decreases, and the tail
action increases with that entry, so $E_h'(u)<0$.
Also $h'(u)<0$ on this positive central interval.

Let $L_-$ be the unforced incoming tail action map used to form the
inverse-tail source difference. It is increasing on its common
action-input interval. The source construction gives exactly
\[
 \mathcal D=h\{L_-^{-1}(E_h(P))-L_-^{-1}(E_h(s))\}.
\]
Put $d=P-s$ and $e=E_h(P)-E_h(s)$. The integral mean-value formula
now gives \eqref{rw:source:positive-unit} with
\begin{equation}\label{rw:source:unit-formula}
 W=h\left[\int_0^1(L_-^{-1})'(E_h(s)+ve)\,\dd v\right]
 \frac{\displaystyle\int_0^1E_h'(s+vd)\,\dd v}
      {\displaystyle\int_0^1h'(s+vd)\,\dd v}>0.
\end{equation}
The nonzero denominators and finite-time analytic transitions give a
positive analytic unit, including $P=s$, on the actual overlap.

At zero parameter, $E_h(u)=h(u)/h$ and the unforced tail is exactly
\[
 T_0(v)=\frac12+v+\frac{e^{-hv}-1}{h},\qquad
 T_0'(v)=1-e^{-hv}>0.
\]
Writing $z_*(h)$ for the base height at $x=\sqrt h$, we obtain
\begin{equation}\label{rw:source:base-unit}
 W(h,0)=\frac1{1-z_*(h)}>0,
 \qquad z_*(h)=s(h/2).
\end{equation}
On a fixed overlap, shrink the original coefficient neighborhood
so that all signed line-breaking segments through zero lie inside
the analytic compact return tube and the independent source domains.
Hadamard division in $\ell$ then defines $\Phi/\ell$ with its
two required energy derivatives. The unforced source uses the same
$(\tau_0,a,t,B)$ division as the preceding lemma. Its first and mixed
columns at zero are therefore the corresponding first and mixed
derivatives of $W\Delta$. The only possible additional mixed terms
are derivatives of $W$ times $\Delta$, $\partial_a\Delta$ or
$\partial_t\Delta$ at zero, all of which vanish by the center
identities and Appendix~\ref{rw:appendix-base-ect}. Thus the mixed
column is $W(h,0)\mathfrak c$; the other two columns are
$W(h,0)\mathfrak a$ and $W(h,0)\mathfrak b$.
Analytic division gives the stated $C^2$ convergence.
\end{proof}

\begin{proof}[Proof of Proposition~\ref{rw:initial-joint-bound}]
The constants are chosen in a fixed order. Appendix
\ref{rw:appendix-base-ect} first fixes $H_{\rm ECT}$, and the
collar gives $h_L:=h(z_L)>H_{\rm ECT}+2$. Choose the small source
coefficient box and then increase its threshold to a fixed
$H_*>h_L+1$, retaining a little more source range to the left of
$H_*$. The compact interval $[h_L,H_*+1]$ and a fixed enlargement
are now fixed. Only then shrink the original coefficient neighborhood.
Make this shrink reflection invariant and small enough that the fixed
overlap $[H_*,H_*+1]$ belongs to all the retained domains: its scaled
coefficients and $|\ell|/s(h)$ are uniformly small, and both core
values are close to one there.

Fix a field and first use the direction $\ell\ge0$.
The source curvature and localization results give one positive-core
source interval $J_s$ starting at $H_*$ and containing that overlap.
Its upper endpoint is determined by the scaled-coefficient and core
conditions, not by any closing zero. It may be infinite. The compact
return on $[h_L,H_*+1]$ and the source return agree on their overlap,
so together they give one actual positive return on the interval
\[
 J=[h_L,H_*]\cup J_s.
\]
Every counted cycle of the first two classes is represented on $J$.
For the source class, Lemma~\ref{rw:geometry:cycle-height} gives
$|\widehat\ell|<3$, and Proposition~\ref{rw:source:localization}
puts it strictly inside the required core inequalities. Hence no
counted zero of those classes is discarded by forming $J$.

Extend $W$ from the source part to a positive smooth function
$\widetilde W$ on $J$, agreeing with $W$ for $h\ge H_*$.
For example, use its already retained slightly larger source overlap,
multiply $\log W$ by a fixed cutoff which vanishes to the left of
that overlap and is one by $H_*$, and exponentiate. Define one
function
\[
 F=\widetilde W\Delta.
\]
Choose a fixed smooth cutoff $\chi_0(h)$ equal to one for
$h\le H_*$ and zero for $h\ge H_*+1$. On the fixed overlap both
coefficient representations are defined. Extend their cutoff products
by their respective formulas outside it, and put
\begin{equation}\label{rw:source:glued-columns}
\begin{aligned}
 f_1&=\chi_0\widetilde W\mathcal A+(1-\chi_0)f,\\
 f_2&=\chi_0\widetilde W\mathcal C+(1-\chi_0)qf,\\
 f_3&=\chi_0\ell\widetilde W\mathcal B-(1-\chi_0)\Phi.
\end{aligned}
\end{equation}
The exact identities \eqref{rw:source:compact-columns} and
\eqref{rw:source:positive-unit} give, throughout $J$,
\begin{equation}\label{rw:source:single-F}
 F=\tau_0 f_1+\sigma_0 f_2+f_3.
\end{equation}
Thus the cutoff has changed a representation of the same function.

We next verify, on the whole interval,
\begin{equation}\label{rw:source:joined-wronskians}
 f_1>0,\qquad W[f_1,f_2]>0,\qquad
 \operatorname{sign}W[f_1,f_2,f_3]=-\operatorname{sign}\ell
 \quad(\ell\ne0).
\end{equation}
On the compact left portion this follows from
Lemma~\ref{rw:source:compact-division} and the identities
$W[wa,wc]=w^2W[a,c]$ and
$W[wa,wc,wb]=w^3W[a,c,b]$ for a common positive multiplier.
On the source right portion, Proposition~\ref{rw:source:curvature}
gives $q_h>0$ and, with $g=\Phi/f$, strict positive $g_{qq}$
when $\ell>0$. Direct differentiation gives
\[
 W[f,qf]=f^2q_h>0,\qquad
 W[f,qf,-\Phi]=-f^3q_h^3g_{qq}<0.
\]
These derivatives are taken at fixed original coefficients, as in
the physical derivative specified in the source construction.

On $[H_*,H_*+1]$, divide the third column in
\eqref{rw:source:glued-columns} by $\ell$ before taking a limit.
By Lemma~\ref{rw:source:common-displacement}, the two triples being
combined have the identical $C^2$ limit
\eqref{rw:source:common-base-triple}. Since the cutoff is fixed,
their combination has that limit as well. Its ordered Wronskians
have strict margins on this fixed compact interval, by
\eqref{rw:ect-wronskian-signs} and $W(h,0)>0$.
One common coefficient shrink therefore proves
\eqref{rw:source:joined-wronskians} also on the overlap. The same
argument for the first two columns applies at $\ell=0$.

Finally, set $r=f_2/f_1$. For $\ell>0$, equation
\eqref{rw:source:single-F} and the Wronskian identities imply
\[
 r'>0,\qquad
 \left(\frac{(F/f_1)'}{r'}\right)'
 =\frac{W[f_1,f_2,f_3]}{f_1^3(r')^2}<0.
\]
The quotient and Rolle argument used in
Proposition~\ref{rw:base-ect} now gives at most two zeros of $F$,
including finite endpoints. Although the cutoff multiplier is only
smooth, it is positive: zero orders up to two agree with those of
$\Delta$, and an order at least three would contradict the strict
displayed derivative. On an unbounded $J$ use each finite selection
of zeros, as in Lemma~\ref{rw:counting:rolle}.

When $\ell=0$, $f_3=0$ and
$F/f_1=\tau_0+\sigma_0r$. If $\sigma_0\ne0$, its one possible
zero is simple; if $\sigma_0=0\ne\tau_0$, it has no zero.
When both vanish, the exact center identities give $F=0$ and no
isolated return. Positive multiplication and the analytic energy
coordinate preserve the Poincar\'e multiplicities.

For an original field with $\ell<0$, apply the proved statement to
its single reflected field. Reflection fixes the central height $s$,
the energy $h$ and the sizes of the four scaled coefficients, so it
preserves the two initial classes. Near each periodic point its
forward return is the inverse of the original forward return, as in
\eqref{rw:geometry:inverse-germ}; that identity preserves the order
of every isolated displacement zero. These local inverse germs
transfer the joint bound to the original field, while the interval
of positive returns remains in the selected reflected direction.
\end{proof}
